\subsection{spin-1}

在 3.4 节中, 我们已经完成了零质量粒子因果场的一般构造。本节将这一框架应用于自旋-1 的零质量粒子——光子。
我们将直接从 Lorentz 群的 $(1,0)\oplus(0,1)$ 表示出发, 构造场强张量 $f_{\mu\nu}$, 
从而得到一个\textbf{不含规范冗余}的自由光子场描述, 之后再讨论矢量势 $A^\mu$ 及其规范对称性的群论起源。

\topictitle{螺旋度与表示的对应}
由 3.4 节的核心结果 (3.93), $(A,B)$ 表示中零质量粒子的螺旋度为 $\lambda = B - A$。
对于自旋-1 粒子 ($\lambda = \pm 1$), 有两种最小的不可约表示：

\begin{center}
\begin{tabular}{c c c}
\hline
表示 & $\lambda$ & 物理含义 \\
\hline
$(1,0)$ & $0-1=-1$ & 反自对偶 (anti-self-dual) \\
$(0,1)$ & $1-0=+1$ & 自对偶 (self-dual) \\
\hline
\end{tabular}
\end{center}
由 3.2 节关于空间反演的讨论, 宇称变换将 $$\mathbb{P}:(A,B) \to (B,A)\to\mathbb{P}: (1,0) \leftrightarrow (0,1)$$
因此, 为了保持空间反演对称性, 必须取直和表示
\begin{equation}
    (1,0) \oplus (0,1)
\end{equation}
这是一个 $3+3=6$ 维表示, 恰好对应四维反对称张量 $f_{\mu\nu} = -f_{\nu\mu}$ 的 6 个独立分量。


\topictitle{自对偶与反自对偶分解}
\textbf{反对称张量的独立分量。}
四维反对称张量 $f_{\mu\nu}=-f_{\nu\mu}$ 有 6 个独立分量, 
恰好可以组织为两组三维矢量。约定
\begin{equation}
    E_i \equiv f_{0i}, \qquad B_k \equiv -\frac{1}{2}\epsilon_{ijk}f_{ij}
    \label{eq:EB-def}
\end{equation}
即 $f_{ij} = -\epsilon_{ijk}B_k$。逆关系 $B_k = \frac{1}{2}\epsilon_{ijk}f_{ij}$ 
(例如 $B_1 = f_{32}$, $B_2 = f_{13}$, $B_3 = f_{21}$)。
将全部分量列出：
\begin{equation}
    f_{\mu\nu} = \begin{pmatrix}
    0 & E_1 & E_2 & E_3 \\
    -E_1 & 0 & -B_3 & B_2 \\
    -E_2 & B_3 & 0 & -B_1 \\
    -E_3 & -B_2 & B_1 & 0
    \end{pmatrix}
    \label{eq:f-matrix}
\end{equation}
\textbf{Hodge 对偶} ：定义$f_{\mu\nu}$ 的 Hodge 对偶张量为
\begin{equation}
    (\star f)_{\mu\nu} \equiv \frac{1}{2}\,\epsilon_{\mu\nu\rho\sigma}\,f^{\rho\sigma}
    \label{eq:hodge-def}
\end{equation}
其中 $\epsilon_{0123}=-1$ ($\epsilon^{0123}=+1$)。
显式计算几个分量：
\begin{enumerate}
    \item $(\star f)_{01}$:
        \[
            \begin{aligned}
            ((\star f))_{0i}
                &= \frac{1}{2}\,\epsilon_{0i\rho\sigma}\,f^{\rho\sigma} = \frac{1}{2}\,\epsilon_{0ijk}\,f^{jk} = \frac{1}{2}\,\epsilon_{ijk}\,f^{jk} = B_i
            \end{aligned}
        \]
    \item   $(\star f)_{23}$:
        \[
                \begin{aligned}
                ((\star f))_{ij}
                &= \frac{1}{2}\,\epsilon_{ij\rho\sigma}\,f^{\rho\sigma} = \frac{1}{2}\bigl(\epsilon_{ij0k}\,f^{0k} + \epsilon_{ijk0}\,f^{k0}\bigr) \\
                &= \frac{1}{2}\bigl(\epsilon_{ij0k}\,f^{0k} - \epsilon_{ij0k}\,f^{0k}\bigr)\,(-1) = \epsilon_{ij0k}\,f^{0k} \\
                &= -\,\epsilon_{ijk}\,f^{0k} = \epsilon_{ijk}\,E_k
            \end{aligned}
        \]
\end{enumerate}
完成全部分量的计算后, 对偶张量具有与 (\ref{eq:f-matrix}) 相同的结构, 
但 $\vec{E}$ 和 $\vec{B}$ 互换并带上符号：
\begin{equation}
    (\star f)_{\mu\nu}: \quad \vec{E} \to \vec{B}, \qquad \vec{B} \to -\vec{E}
    \label{eq:hodge-EB}
\end{equation}
矩阵形式为
\begin{equation}
    (\star f)_{\mu\nu} = \begin{pmatrix}
    0 & B_1 & B_2 & B_3 \\
    -B_1 & 0 & E_3 & -E_2 \\
    -B_2 & -E_3 & 0 & E_1 \\
    -B_3 & E_2 & -E_1 & 0
    \end{pmatrix}
\end{equation}
\textbf{Hodge 对偶的平方}: 再作用一次 Hodge 对偶，有
\begin{equation}
    \star^2 f_{\mu\nu}:\quad 
    \vec{E} \xrightarrow{\star} \vec{B} \xrightarrow{\star} -\vec{E}, \qquad
    \vec{B} \xrightarrow{\star} -\vec{E} \xrightarrow{\star} -\vec{B}
\end{equation}
因此
\begin{equation}
    \star^2 f_{\mu\nu} = - f_{\mu\nu},
\end{equation}
即
\begin{equation}
    \star^2 f = -f 
    \qquad 
    \bigl(\text{闵氏时空中 Hodge 对偶的平方为 } -1\bigr).
    \label{eq:hodge-squared}
\end{equation}
也就是\[
\star^2 f_{\mu\nu}
= \frac{1}{2}\epsilon_{\mu\nu\rho\sigma}
\left(\frac{1}{2}\epsilon^{\rho\sigma\alpha\beta} f_{\alpha\beta}\right)
= - f_{\mu\nu}
\]
\textbf{投影算符与分解.}
由于 $\star^2 = -1$，算符 $i\star$ 满足 $(i\star)^2 = 1$，
其本征值为 $\pm 1$。定义投影算符
\begin{equation}
\mathbb{P}^\pm_{\mu\nu}{}^{\rho\sigma}
= \frac{1}{2}\left(
\delta^\rho_\mu\delta^\sigma_\nu
- \delta^\rho_\nu\delta^\sigma_\mu
\pm \frac{i}{2}\,\epsilon_{\mu\nu}{}^{\rho\sigma}
\right),
\end{equation}
它们满足 $\mathbb{P}^+ + \mathbb{P}^- = \mathbb{I}$，
$\mathbb{P}^+\mathbb{P}^- = 0$，$(\mathbb{P}^\pm)^2 = \mathbb{P}^\pm$。
作用在反对称张量 $f_{\mu\nu}$ 上，得到分解
\begin{equation}
f^\pm_{\mu\nu}
\equiv \mathbb{P}^\pm_{\mu\nu}{}^{\rho\sigma} f_{\rho\sigma}
= \frac{1}{2}\left(f_{\mu\nu} \pm i(\star f)_{\mu\nu}\right),
\end{equation}
从而
\begin{equation}
f_{\mu\nu} = f^+_{\mu\nu} + f^-_{\mu\nu}, \qquad
i\star f^\pm_{\mu\nu} = \pm f^\pm_{\mu\nu}.
\end{equation}
这给出了反对称张量在复化空间中的分解
\[
\boxed{\Lambda^2 \otimes \mathbb{C} = \Lambda^2_+ \oplus \Lambda^2_-,}
\]
对应 Lorentz 群的 $(1,0)$ 与 $(0,1)$ 不可约表示。在三维矢量语言下，该分解对应
\begin{equation}
\boldsymbol{f}^{\,\pm} = \frac{1}{2}(\boldsymbol{E} \pm i\boldsymbol{B}),
\end{equation}
即 Riemann--Silberstein 复矢量，它们分别携带 $(1,0)$ 与 $(0,1)$ 表示，
并满足 $\boldsymbol{f}^{\,-} = (\boldsymbol{f}^{\,+})^*$。
\topictitle{6 维表示空间的不可约分解}
上面的 Riemann--Silberstein 分解 $\boldsymbol{f}^{\,\pm}=\frac{1}{2}(\boldsymbol{E}\pm i\boldsymbol{B})$
是一个代数定义。现在要证明它确实对应 Lorentz 群的 $(0,1)$ 与 $(1,0)$ 不可约表示。
方法是：直接从电磁场的 boost 变换规则出发, 证明两组复矢量各自封闭,
读出每个子空间上的旋转与 boost 生成元, 最后用
$\boldsymbol{A}=\frac{1}{2}(\boldsymbol{J}-i\boldsymbol{K})$,
$\boldsymbol{B}=\frac{1}{2}(\boldsymbol{J}+i\boldsymbol{K})$ 认出 $(A,B)$。

\textbf{电磁场在 boost 下的变换。}
沿 $z$ 方向快度为 $\phi$ 的标准 boost 下,
$f'_{\mu\nu}=\Lambda^\rho_{\ \mu}\Lambda^\sigma_{\ \nu}f_{\rho\sigma}$ 给出
\begin{subequations}\label{eq:boost-EB}
\begin{align}
    &E'_1 = \cosh\phi\,E_1 - \sinh\phi\,B_2, \quad 
     E'_2 = \cosh\phi\,E_2 + \sinh\phi\,B_1, \quad 
     E'_3 = E_3 \label{eq:boost-E}\\
    &B'_1 = \sinh\phi\,E_2 + \cosh\phi\,B_1, \quad  
     B'_2 = -\sinh\phi\,E_1 + \cosh\phi\,B_2, \quad  
     B'_3 = B_3 \label{eq:boost-B}
\end{align}
\end{subequations}
关键观察：boost \textbf{混合} $\vec{E}$ 和 $\vec{B}$
(例如 $E_1$ 与 $B_2$ 耦合), 而旋转把 $\vec{E}$ 和 $\vec{B}$ 各自内部旋转。
因此 6 维空间 $(\vec{E},\vec{B})$ 在 Lorentz 群下不是不可约的——
应该能拆成两个 3 维不可约子空间。

\textbf{验证 $\boldsymbol{f}^+$ 自洽封闭。}
计算 $f'^+_1 = \frac{1}{2}(E'_1 + iB'_1)$：
\begin{align}
    f'^+_1 &= \frac{1}{2}\bigl[(\cosh\phi\,E_1 - \sinh\phi\,B_2) 
    + i(\sinh\phi\,E_2 + \cosh\phi\,B_1)\bigr] \nonumber\\
    &= \frac{1}{2}\bigl[\cosh\phi\,(E_1+iB_1) + \sinh\phi\,(-B_2+iE_2)\bigr] \nonumber\\
    &= \cosh\phi\,f^+_1 + i\sinh\phi\,f^+_2
\end{align}
其中第二步利用了 $-B_2+iE_2 = i(E_2+iB_2) = 2if^+_2$。类似地：
\begin{equation}
    f'^+_2 = i\sinh\phi\,f^+_1 + \cosh\phi\,f^+_2, \qquad f'^+_3 = f^+_3
\end{equation}
$\vec{f}^{\,+}$ 的三个分量\textbf{只和自己混合}, 完全不涉及 $\vec{f}^{\,-}$。
对 $\vec{f}^{\,-}$ 做同样计算 ($i\to -i$), 也自洽封闭。
这就证明了 6 维空间在 boost 下分裂为两个 3 维不变子空间。

\textbf{旋转生成元。}
旋转不区分 $\vec{E}$ 和 $\vec{B}$ (两者都是三维矢量, 变换规则相同),
因此 $\vec{f}^{\,\pm}$ 在旋转下的生成元都是标准的 $j=1$ 角动量矩阵 $J_i^{(1)}$,
在 Cartesian 基 $(f^{(\pm)}_1, f^{(\pm)}_2, f^{(\pm)}_3)$ 下为
\begin{equation}
    (J^{(1)}_k)_{ij} = -i\epsilon_{kij}
\end{equation}
显式写出 $J^{(1)}_3$：
\begin{equation}
    J^{(1)}_3 = \begin{pmatrix} 0 & -i & 0 \\ i & 0 & 0 \\ 0 & 0 & 0 \end{pmatrix}
    \qquad\Longrightarrow\qquad
    J^{(1)}_3 \begin{pmatrix} f_1 \\ f_2 \\ f_3 \end{pmatrix}
    = \begin{pmatrix} -if_2 \\ if_1 \\ 0 \end{pmatrix}
\end{equation}

\textbf{引入球面基与升降算符。}
Cartesian 基 $(f_1,f_2,f_3)$ 不是 $J_3$ 的本征态。
为了对角化 $J_3$, 定义球面基 (spherical basis)：
\begin{equation}
    f_+ \equiv -\frac{1}{\sqrt{2}}(f_1 + if_2), \qquad
    f_0 \equiv f_3, \qquad
    f_- \equiv \frac{1}{\sqrt{2}}(f_1 - if_2)
    \label{eq:spherical-basis-def}
\end{equation}
逆关系为
\begin{equation}
    f_1 = -\frac{1}{\sqrt{2}}(f_+ - f_-), \qquad
    f_2 = -\frac{i}{\sqrt{2}}(f_+ + f_-), \qquad
    f_3 = f_0
    \label{eq:spherical-basis-inv}
\end{equation}
验证 $J_3$ 在球面基下是对角的：
\begin{align}
    J_3 f_+ &= -\frac{1}{\sqrt{2}}(J_3 f_1 + iJ_3 f_2)
    = -\frac{1}{\sqrt{2}}(-if_2 + i\cdot if_1) \nonumber\\
    &= -\frac{1}{\sqrt{2}}(-if_2 - f_1)
    = -\frac{1}{\sqrt{2}}(f_1 + if_2) \cdot (+1) = +1\cdot f_+
\end{align}
类似地 $J_3 f_0 = 0$, $J_3 f_- = -1\cdot f_-$。因此
\begin{equation}
    \boxed{J_3\,f_m = m\,f_m, \qquad m = +1,\,0,\,-1}
    \label{eq:J3-diag}
\end{equation}
球面基正是 $J_3$ 的\textbf{权基} (weight basis), $m$ 是权值 (weight)。

定义升降算符 $J_\pm = J_1 \pm iJ_2$, 它们满足标准 $\mathfrak{su}(2)$ 代数
$[J_3, J_\pm] = \pm J_\pm$, $[J_+, J_-] = 2J_3$, 以及
\begin{equation}
    J_\pm\,f_m = \sqrt{(1\mp m)(1\pm m+1)}\;f_{m\pm 1}
    \label{eq:J-pm-action}
\end{equation}
显式地：
\begin{equation}
    J_+\,f_- = \sqrt{2}\,f_0, \quad J_+\,f_0 = \sqrt{2}\,f_+, \quad J_+\,f_+ = 0
\end{equation}
以及 $J_-$ 是 $J_+$ 的厄米共轭, 方向相反。
$\{f_+, f_0, f_-\}$ 构成 $j=1$ 表示的标准权基, $J_+$ 和 $J_-$ 在基底之间阶梯式移动。

\textbf{Boost 生成元在球面基下的作用。}
从 Step 2 的变换结果, 转到球面基。利用 (\ref{eq:spherical-basis-def}):
\begin{align}
    f'^{(\pm)}_+ &= -\frac{1}{\sqrt{2}}(f'^{(\pm)}_1 + if'^{(\pm)}_2)
\end{align}
对 $\vec{f}^{\,+}$, 代入 $f'^+_1 = \cosh\phi\,f^+_1 + i\sinh\phi\,f^+_2$
和 $f'^+_2 = i\sinh\phi\,f^+_1 + \cosh\phi\,f^+_2$：
\begin{align}
    f'^+_+ &= -\frac{1}{\sqrt{2}}\bigl[(\cosh\phi + \sinh\phi)(f^+_1 + if^+_2)\bigr]
    = e^{+\phi}\,f^+_+ \\
    f'^+_0 &= f^+_0 \\
    f'^+_- &= -\frac{1}{\sqrt{2}}\bigl[(\cosh\phi - \sinh\phi)(f^+_1 - if^+_2)\bigr] \cdot(-1)
    = e^{-\phi}\,f^+_-
\end{align}
统一写为
\begin{equation}
    f'^+_m = e^{+m\phi}\,f^+_m
    \label{eq:fplus-boost}
\end{equation}
对 $\vec{f}^{\,-}$ ($i\to -i$)：
\begin{equation}
    f'^-_m = e^{-m\phi}\,f^-_m
    \label{eq:fminus-boost}
\end{equation}

\textbf{从变换因子读出生成元。}
有限 boost 算符 $e^{-i\phi K_3}$ 在球面基上的作用是 $f_m \to e^{-i\phi K_3}f_m$。
比较 (\ref{eq:fplus-boost}), 要求 $e^{-i\phi K_3^{(+)}}$ 在 $f^+_m$ 上给出因子 $e^{+m\phi}$,
即 $-iK_3^{(+)} = +J_3^{(1)}$\,\footnote{%
因为 $J_3^{(1)}$ 在权基上作用为 $J_3 f_m = mf_m$,
所以 $e^{+\phi J_3}f_m = e^{+m\phi}f_m$。}, 所以
\begin{equation}
    K_3^{(\text{on }\vec{f}^+)} = -iJ_3^{(1)}, \qquad
    K_3^{(\text{on }\vec{f}^-)} = +iJ_3^{(1)}
\end{equation}
推广到三个方向\footnote{%
对沿 $x$ 和 $y$ 方向的 boost 做同样计算, 或利用 $[J_i, K_j]=i\epsilon_{ijk}K_k$
的代数一致性, 均可得到此推广。}：
\begin{equation}
    \vec{K}^{(\text{on }\vec{f}^+)} = -i\vec{J}^{(1)}, \qquad
    \vec{K}^{(\text{on }\vec{f}^-)} = +i\vec{J}^{(1)}
    \label{eq:K-on-fpm}
\end{equation}

\textbf{认出 $(A,B)$ 表示。}
利用 $\boldsymbol{A}=\frac{1}{2}(\boldsymbol{J}-i\boldsymbol{K})$,
$\boldsymbol{B}=\frac{1}{2}(\boldsymbol{J}+i\boldsymbol{K})$,
代入每个子空间上的旋转和 boost 生成元。

对 $\vec{f}^{\,+}$ ($\vec{J}=\vec{J}^{(1)}$, $\vec{K}=-i\vec{J}^{(1)}$)：
\begin{align}
    A_i &= \tfrac{1}{2}\bigl(J_i^{(1)} - i(-iJ_i^{(1)})\bigr) 
    = \tfrac{1}{2}(J_i^{(1)} - J_i^{(1)}) = 0 \\
    B_i &= \tfrac{1}{2}\bigl(J_i^{(1)} + i(-iJ_i^{(1)})\bigr) 
    = \tfrac{1}{2}(J_i^{(1)} + J_i^{(1)}) = J_i^{(1)}
\end{align}
因此 $A=0$, $B=1$: $\vec{f}^{\,+}$ 按 $(0,1)$ 变换。$\checkmark$

对 $\vec{f}^{\,-}$ ($\vec{J}=\vec{J}^{(1)}$, $\vec{K}=+i\vec{J}^{(1)}$)：
\begin{align}
    A_i &= \tfrac{1}{2}\bigl(J_i^{(1)} - i(+iJ_i^{(1)})\bigr) 
    = \tfrac{1}{2}(J_i^{(1)} + J_i^{(1)}) = J_i^{(1)} \\
    B_i &= \tfrac{1}{2}\bigl(J_i^{(1)} + i(+iJ_i^{(1)})\bigr) 
    = \tfrac{1}{2}(J_i^{(1)} - J_i^{(1)}) = 0
\end{align}
因此 $A=1$, $B=0$: $\vec{f}^{\,-}$ 按 $(1,0)$ 变换。$\checkmark$

\textbf{与 $(A,B)$ 表示的 boost 矩阵一致性检验。}
在 $(A,B)$ 表示中 (参考 (3.64)),
$\mathscr{D}(L(p)) = \exp(-\phi\cdot\boldsymbol{A})\otimes\exp(+\phi\cdot\boldsymbol{B})$。
\begin{itemize}
    \item $(0,1)$: $A=0$, $B=1$, boost 矩阵 $e^{+\phi B_3}$, 
    本征值 $b=+1,0,-1$, 变换因子 $e^{+b\phi}$:
    \begin{equation}
        e^{+\phi},\;1,\;e^{-\phi} \quad\longleftrightarrow\quad 
        f'^+_+ = e^{+\phi}f^+_+,\;f'^+_0 = f^+_0,\;f'^+_- = e^{-\phi}f^+_-
        \quad\checkmark
    \end{equation}
    
    \item $(1,0)$: $A=1$, $B=0$, boost 矩阵 $e^{-\phi A_3}$, 
    本征值 $a=+1,0,-1$, 变换因子 $e^{-a\phi}$:
    \begin{equation}
        e^{-\phi},\;1,\;e^{+\phi} \quad\longleftrightarrow\quad 
        f'^-_+ = e^{-\phi}f^-_+,\;f'^-_0 = f^-_0,\;f'^-_- = e^{+\phi}f^-_-
        \quad\checkmark
    \end{equation}
\end{itemize}

由此确认：
\begin{equation}
    \boxed{
    \underbrace{f^+_{\mu\nu}}_{(0,1),\;\lambda=+1} \oplus 
    \underbrace{f^-_{\mu\nu}}_{(1,0),\;\lambda=-1} = f_{\mu\nu}
    }
\end{equation}

整条验证路径的逻辑链为：
\begin{equation*}
    \underbrace{f'_{\mu\nu}=\Lambda\Lambda f}_{\text{boost 变换规则}}
    \xrightarrow{\;\vec{f}^{\,\pm}=\frac{1}{2}(\vec{E}\pm i\vec{B})\;}
    \underbrace{\text{各自封闭}}_{\text{两个不变子空间}}
    \xrightarrow{\;\text{球面基}\;}
    \underbrace{e^{\pm m\phi}}_{\text{变换因子}}
    \xrightarrow{\;\vec{A},\vec{B}\;}
    \underbrace{(0,1)\oplus(1,0)}_{\text{认出表示}}
\end{equation*}


\topictitle{极化张量}

由 3.4 节的一般结果, 零质量场的系数函数由标准动量处的值经标准 boost 确定。
对于场强张量, 需要构造反对称极化张量 $e_{\mu\nu}(\mathbf{k},\lambda)$,
它编码了螺旋度为 $\lambda$ 的单光子态的电磁场构型。

\textbf{约束条件。}$e_{\mu\nu}$ 必须满足：
\begin{enumerate}
    \item 反对称性: $e_{\mu\nu}=-e_{\nu\mu}$;
    \item 横向性: $k^\mu e_{\mu\nu}=0$ (对应运动方程 $\partial^\mu f_{\mu\nu}=0$);
    \item 确定的对偶性: $\lambda=+1$ 时纯自对偶, $\lambda=-1$ 时纯反自对偶
    (来自 ISO(2) 约束 (3.91))。
\end{enumerate}

\textbf{一般形式与唯一性。}
在标准动量 $k^\mu=(\kappa,0,0,\kappa)$ 处,
能用来构造反对称张量的材料只有 $k_\mu$ 和横向圆极化矢量
$\epsilon_\mu(\mathbf{k},\lambda)$\,\footnote{%
圆极化矢量约定: $\epsilon^\mu(+1)=\frac{1}{\sqrt{2}}(0,-1,-i,0)$,
$\epsilon^\mu(-1)=\frac{1}{\sqrt{2}}(0,+1,-i,0)$;
满足 $k\cdot\epsilon=0$, $\epsilon\cdot\epsilon^*=-1$。}。
最一般的反对称组合为\footnote{%
$k_\mu k_\nu$ 和 $\epsilon_\mu\epsilon_\nu$ 反对称化后为零, 故上式已是最一般形式。
两项均自动满足横向性。}
\begin{equation}
    e_{\mu\nu} = \alpha\,\underbrace{(k_\mu\epsilon_\nu - k_\nu\epsilon_\mu)}_{T^{(1)}_{\mu\nu}}
    + \beta\,\underbrace{\epsilon_{\mu\nu\rho\sigma}\,k^\rho\epsilon^\sigma}_{T^{(2)}_{\mu\nu}}
    \label{eq:general-ansatz}
\end{equation}
圆极化矢量 $\epsilon_\mu(\lambda)$ 本身已编码了手征性——在标准动量处可以直接计算验证\footnote{%
以 $\lambda=+1$ 为例: $T^{(2)}_{01} = \epsilon_{0132}k^3\epsilon^2(+1)
= (+1)\kappa\frac{-i}{\sqrt{2}} = \frac{-i\kappa}{\sqrt{2}}$;
$T^{(1)}_{01} = k_0\epsilon_1(+1) = \kappa\frac{1}{\sqrt{2}}$;
比值 $T^{(2)}_{01}/T^{(1)}_{01}=-i$。验证其余分量得到同样比值。}：
\begin{equation}
    \epsilon_{\mu\nu\rho\sigma}\,k^\rho\epsilon^\sigma(\lambda)
    = -i\lambda\,(k_\mu\epsilon_\nu(\lambda) - k_\nu\epsilon_\mu(\lambda))
    \label{eq:duality-auto}
\end{equation}
代入 (\ref{eq:general-ansatz}), 两项合并为同一结构：
$e_{\mu\nu} = (\alpha - i\lambda\beta)\,T^{(1)}$。
因此对偶性被圆极化矢量自动编码, 只剩一个自由常数。
选取 $\alpha=i$, $\beta=0$, 得到
\begin{equation}
    \boxed{e_{\mu\nu}(\mathbf{k},\lambda) =
    i\bigl(k_\mu\,\epsilon_\nu(\mathbf{k},\lambda) - k_\nu\,\epsilon_\mu(\mathbf{k},\lambda)\bigr)}
    \label{eq:pol-tensor}
\end{equation}
尽管形式上与 ``$\partial_\mu A_\nu - \partial_\nu A_\mu$" 在动量空间中相同,
但这里 $A_\mu$ 尚未被引入——(\ref{eq:pol-tensor})
完全是 $(1,0)\oplus(0,1)$ 表示论和 ISO(2) 约束的产物。

\textbf{与 3.4 节系数函数 $v_{ab}$ 的对应。}
极化张量是 3.4 节一般系数函数 $v_{ab}(\mathbf{k},\lambda)$
在反对称张量基底中的表达。推导分为四步。

\textit{Step 1: 球面基与 Cartesian 基的转换。}
$(0,1)$ 扇区的表示空间是三维的 ($B=1$),
基底 $|b\rangle$ ($b=+1,0,-1$) 与 Cartesian 基 $\hat{e}_i$ 的关系
正是上面 (\ref{eq:spherical-basis-def}) 定义的球面基：
\begin{equation}
    |b=+1\rangle = -\frac{1}{\sqrt{2}}(\hat{e}_1 + i\hat{e}_2), \qquad
    |b=0\rangle = \hat{e}_3, \qquad
    |b=-1\rangle = \frac{1}{\sqrt{2}}(\hat{e}_1 - i\hat{e}_2)
\end{equation}
$(1,0)$ 扇区 ($A=1$) 类似, 用 $|a\rangle$ ($a=+1,0,-1$) 标记。

\textit{Step 2: ISO(2) 约束筛选分量。}
由 (3.91), 零质量条件将允许的 $(a,b)$ 值锁定为 $a=-A$, $b=+B$:
\begin{itemize}
    \item $(0,1)$ 扇区, $\lambda=+1$: $A=0,B=1 \Rightarrow a=0,\;b=+1$。
    CG 系数 $\langle 0,0;\,1,+1\,|\,1,+1\rangle = 1$, 仅 $b=+1$ 分量非零。
    \item $(1,0)$ 扇区, $\lambda=-1$: $A=1,B=0 \Rightarrow a=-1,\;b=0$。
    CG 系数 $\langle 1,-1;\,0,0\,|\,1,-1\rangle = 1$, 仅 $a=-1$ 分量非零。
\end{itemize}

\textit{Step 3: 从球面基到场强分量。}
$(0,1)$ 扇区的 Cartesian 分量对应 $\vec{f}^{\,+}=\frac{1}{2}(\vec{E}+i\vec{B})$。
仅 $b=+1$ 非零, 利用球面基：
\begin{equation}
    \vec{f}^{\,+} \propto |b=+1\rangle = -\frac{1}{\sqrt{2}}(\hat{e}_1 + i\hat{e}_2)
\end{equation}
由纯自对偶条件 $\vec{E}-i\vec{B}=0$ (即 $\vec{B}=-i\vec{E}$),
$E_i+iB_i = 2E_i$, 所以
\begin{equation}
    \vec{E} \propto -\frac{1}{\sqrt{2}}(\hat{e}_1 + i\hat{e}_2)
    = \frac{1}{\sqrt{2}}(-1,\,-i,\,0)
\end{equation}
这与圆极化矢量 $\epsilon^\mu(+1)=\frac{1}{\sqrt{2}}(0,-1,-i,0)$
的空间部分精确一致。

\textit{Step 4: 组装张量指标。}
$e_{0i}\sim E_i$ 和 $e_{ij}\sim -\epsilon_{ijk}B_k$ 的关系,
加上在壳条件 $k^2=0$ 对 $e_{3i}$ 分量的约束\footnote{%
在标准动量 $k^\mu=(\kappa,0,0,\kappa)$ 处, $k_0=k_3=\kappa$,
所以 $e_{0i}=i\kappa\epsilon_i$ 而 $e_{3i}=-i\kappa\epsilon_i$,
即 $e_{3i}=-e_{0i}$——这正是 $k^2=0$ 的推论。},
最终将全部 $v_{ab}$ 分量组装成 $e_{\mu\nu}=i(k_\mu\epsilon_\nu - k_\nu\epsilon_\mu)$。

总结：
\begin{equation*}
    \underbrace{v_{ab}(\mathbf{k},\lambda)}_{\text{3.4 节, 球面基}}
    \xrightarrow{\;\text{ISO(2) 筛选}\;}
    \underbrace{b=+1 \text{ 或 } a=-1}_{\text{仅一个分量非零}}
    \xrightarrow{\;\text{球面基}\to\text{Cartesian}\;}
    \underbrace{\vec{E},\,\vec{B}}_{\text{场强分量}}
    \xrightarrow{\;\text{组装}\;}
    \underbrace{e_{\mu\nu}}_{\text{反对称张量}}
\end{equation*}

\textbf{验证自对偶性。}以 $\lambda=+1$ 为例。取 $k^\mu=(\kappa,0,0,\kappa)$,
$\epsilon_\mu = \eta_{\mu\nu}\epsilon^\nu = \frac{1}{\sqrt{2}}(0,1,i,0)$。
由 (\ref{eq:pol-tensor}) 计算非零分量：
\begin{align}
    e_{01} &= i\,k_0\epsilon_1 = \frac{i\kappa}{\sqrt{2}}, &\quad
    e_{02} &= i\,k_0\epsilon_2 = -\frac{\kappa}{\sqrt{2}} \\
    e_{31} &= -i\,k_3\epsilon_1 = -\frac{i\kappa}{\sqrt{2}}, &\quad
    e_{32} &= -i\,k_3\epsilon_2 = \frac{\kappa}{\sqrt{2}}
\end{align}
读出 $\vec{E}=\frac{\kappa}{\sqrt{2}}(i,-1,0)$,
$\vec{B}=\frac{\kappa}{\sqrt{2}}(1,i,0)$。验证：
\begin{equation}
    \vec{E}+i\vec{B} = \sqrt{2}\,\kappa\,(i,-1,0) \neq 0, \qquad
    \vec{E}-i\vec{B} = \vec{0} \quad\checkmark
\end{equation}
因此 $\lambda=+1$ 对应纯自对偶 $(0,1)$。类似计算可验证 $\lambda=-1$ 对应纯反自对偶 $(1,0)$。

物理图景：$\lambda=+1$ 的光子, $\vec{E}$ 和 $\vec{B}$ 在横向平面内以右旋方式旋转,
且严格满足 $\vec{B}=-i\vec{E}$; $\lambda=-1$ 则左旋, $\vec{B}=+i\vec{E}$。


\topictitle{场展开}

光子是自身的反粒子 ($a^{c\dagger}=a^\dagger$),
且自旋-统计定理要求整数自旋粒子满足对易关系 (玻色子)。
由 3.4 节的一般结果 (3.94), $(-1)^{2B}=1$ 对两个扇区都成立\footnote{%
$(1,0)$: $B=0$, $(-1)^{2B}=1$; $(0,1)$: $B=1$, $(-1)^{2B}=1$。};
产生部分系数由 (3.100) 给出 $u_{\mu\nu}=v_{\mu\nu}^*$, 即 $e^*_{\mu\nu}$。场展开为：
\begin{equation}
    \boxed{
    f_{\mu\nu}(x) = \frac{1}{(2\pi)^{3/2}} \sum_{\lambda=\pm 1} \int \frac{d^3k}{\sqrt{2k^0}}
    \left[e_{\mu\nu}(\mathbf{k},\lambda)\,a(\mathbf{k},\lambda)\,e^{ik\cdot x}
    + e^*_{\mu\nu}(\mathbf{k},\lambda)\,a^\dagger(\mathbf{k},\lambda)\,e^{-ik\cdot x}\right]
    }
    \label{eq:f-expansion}
\end{equation}
其中 $k^0=|\mathbf{k}|$, 产生湮灭算符满足：
\begin{align}
    [a(\mathbf{k},\lambda),\,a^\dagger(\mathbf{k}',\lambda')] &= \delta^3(\mathbf{k}-\mathbf{k}')\,
    \delta_{\lambda\lambda'} \label{eq:photon-commutation} \\
    [a(\mathbf{k},\lambda),\,a(\mathbf{k}',\lambda')] &=
    [a^\dagger(\mathbf{k},\lambda),\,a^\dagger(\mathbf{k}',\lambda')] = 0
\end{align}

逐项解读：
\begin{itemize}
    \item $a(\mathbf{k},\lambda)\,e^{ik\cdot x}$ (正频): 湮灭动量 $\mathbf{k}$、螺旋度 $\lambda$ 的光子;
    $a^\dagger(\mathbf{k},\lambda)\,e^{-ik\cdot x}$ (负频): 产生同一光子。
    \item $e_{\mu\nu}$, $e^*_{\mu\nu}$: 极化张量, 编码该模式的 $\vec{E}$-$\vec{B}$ 构型。
    \item $1/\sqrt{2k^0}$: Lorentz 不变归一化, 使 $d^3k/(2k^0)$ 构成不变测度。
\end{itemize}
正频与负频的同时出现是微观因果性的要求 (3.1 节 (3.22))。
由于光子是自身的反粒子, 不需引入额外的反粒子算符。

\textbf{厄米性。}对 (\ref{eq:f-expansion}) 取厄米共轭,
$a\leftrightarrow a^\dagger$, $e_{\mu\nu}\leftrightarrow e^*_{\mu\nu}$,
$e^{ikx}\leftrightarrow e^{-ikx}$, 恰好交换两项, 因此
$f_{\mu\nu}^\dagger(x) = f_{\mu\nu}(x)$。
经典极限中, 这对应 $\vec{E}$, $\vec{B}$ 是实矢量场。

\textbf{经典极限。}对相干态 $|\alpha\rangle$
($a(\mathbf{k},\lambda)|\alpha\rangle = \alpha(\mathbf{k},\lambda)|\alpha\rangle$) 取期望值,
得到经典场强张量的 Fourier 展开。对于单色平面波, 期望值化为
$\langle f_{\mu\nu}\rangle \propto e_{\mu\nu}(\mathbf{k}_0,\lambda)\,e^{ik_0\cdot x} + \text{c.c.}$
——传播方向 $\hat{k}_0$、圆极化 $\lambda$ 的经典平面电磁波, 极化张量的命名由此而来。

\textbf{极化求和。}利用 $e_{\mu\nu}=i(k_\mu\epsilon_\nu - k_\nu\epsilon_\mu)$ 展开：
\begin{equation}
    \sum_{\lambda=\pm 1} e_{\mu\nu}(\mathbf{k},\lambda)\,e^*_{\rho\sigma}(\mathbf{k},\lambda)
    = k_\mu k_\rho\,P_{\nu\sigma} - k_\mu k_\sigma\,P_{\nu\rho}
    - k_\nu k_\rho\,P_{\mu\sigma} + k_\nu k_\sigma\,P_{\mu\rho}
\end{equation}
其中 $P_{\mu\nu}=\sum_\lambda \epsilon_\mu\epsilon^*_\nu$。
对于耦合到守恒流 $j^\mu$ ($k_\mu j^\mu=0$) 的过程,
Ward 恒等式允许 $P_{\mu\nu}\to -\eta_{\mu\nu}$。

\textbf{推广到非阿贝尔规范场。}
上述构造对\textbf{任何}零质量自旋-1 粒子都成立, 区别仅在于内部对称性。
对于规范群 $G$ 的规范玻色子, 场强张量携带伴随表示指标 $a=1,\ldots,\dim G$,
极化张量 $e_{\mu\nu}$ 与色指标无关——Lorentz 结构和内部对称性完全解耦。
对易关系推广为
$[a^a(\mathbf{k},\lambda),\,a^{b\dagger}(\mathbf{k}',\lambda')]
= \delta^{ab}\,\delta^3(\mathbf{k}-\mathbf{k}')\,\delta_{\lambda\lambda'}$。
在完整的非阿贝尔理论中, 场强张量包含非线性项
$F^a_{\mu\nu} = \partial_\mu A^a_\nu - \partial_\nu A^a_\mu + gf^{abc}A^b_\mu A^c_\nu$,
自由场展开仅对应其线性化部分。

这个场展开\textbf{不包含任何规范冗余}——
它直接描述了光子 (或胶子) 的两个物理自由度 $\lambda=\pm 1$。


\topictitle{自由 Maxwell 方程}

展开式 (\ref{eq:f-expansion}) 中, 每个平面波模式携带在壳动量 $k^\mu$ ($k^2=0$, $k^0>0$)
和横向极化矢量 ($k\cdot\epsilon=0$)。
这些条件在坐标空间中翻译为：

\begin{center}
\begin{tabular}{l l l}
\hline
方程 & 动量空间等价 & 来源 \\
\hline
$\Box\,f_{\mu\nu}=0$ & $k^2=0$ & 无质量在壳条件 \\[3pt]
$\partial^\mu f_{\mu\nu}=0$ & $k^\mu e_{\mu\nu}=0$ & 横向性 $k\cdot\epsilon=0$ \\[3pt]
$\partial_{[\rho}f_{\mu\nu]}=0$ & $k_{[\rho}\,e_{\mu\nu]}=0$ &
$e_{\mu\nu}$ 的因式分解结构 \\
\hline
\end{tabular}
\end{center}

\textbf{验证源自方程。}$k^\mu e_{\mu\nu} = i(k^2\epsilon_\nu - k_\nu\,k\cdot\epsilon) = 0$\;$\checkmark$

\textbf{验证 Bianchi 恒等式。}
$k_\rho\,e_{\mu\nu}+k_\mu\,e_{\nu\rho}+k_\nu\,e_{\rho\mu}=0$:
\begin{align}
    &k_\rho(k_\mu\epsilon_\nu - k_\nu\epsilon_\mu)
    + k_\mu(k_\nu\epsilon_\rho - k_\rho\epsilon_\nu)
    + k_\nu(k_\rho\epsilon_\mu - k_\mu\epsilon_\rho) \nonumber\\
    &= \cancel{k_\rho k_\mu\epsilon_\nu} - \bcancel{k_\rho k_\nu\epsilon_\mu}
    + \xcancel{k_\mu k_\nu\epsilon_\rho} - \cancel{k_\mu k_\rho\epsilon_\nu}
    + \bcancel{k_\nu k_\rho\epsilon_\mu} - \xcancel{k_\nu k_\mu\epsilon_\rho}
    = 0 \quad\checkmark
\end{align}

物理意义：$\partial^\mu f_{\mu\nu}=0$ 对应 Maxwell 含源组 (真空):
$\nabla\cdot\vec{E}=0$, $\nabla\times\vec{B}-\partial_t\vec{E}=0$;
$\partial_{[\rho}f_{\mu\nu]}=0$ 对应无源组 (Bianchi):
$\nabla\cdot\vec{B}=0$, $\nabla\times\vec{E}+\partial_t\vec{B}=0$。
后者的成立意味着局域存在矢量势 $A^\mu$ 使得 $f_{\mu\nu}=\partial_\mu A_\nu - \partial_\nu A_\mu$
——但这是数学推论, 非出发点。


\topictitle{矢量势与规范对称性的群论起源}

描述光子的另一种方式是使用四矢量场 $A^\mu$, 按 $(\frac{1}{2},\frac{1}{2})$ 表示变换。
角动量加法 $\frac{1}{2}\otimes\frac{1}{2}=0\oplus 1$ 包含 $j=0$ 和 $j=1$ 两个扇区,
共 4 个分量, 但光子只有 $\lambda=\pm 1$ 两个物理自由度,
多出的分量正是规范自由度的来源。

理解规范冗余的关键在于 Wigner 小群 $ISO(2)$。
``平移"生成元 $\Pi_1,\Pi_2$ 在物理态上本征值为零 (排除连续自旋),
但在四矢量表示中, 小群变换作用在 $\epsilon^\mu$ 上时产生额外项：
\begin{equation}
    W^\mu_{\ \nu}\,\epsilon^\nu(\mathbf{k},\lambda) =
    e^{i\lambda\theta}\,\epsilon^\mu(\mathbf{k},\lambda) + \Omega(\theta,\alpha,\beta)\,k^\mu
\end{equation}
第二项 $\propto k^\mu$ 在坐标空间中对应 $A_\mu(x)\to A_\mu(x)+\partial_\mu\Omega(x)$
——这正是规范变换。

\begin{center}
\fbox{\parbox{0.88\textwidth}{
规范对称性不是额外施加的, 而是 Poincar\'{e} 群在零质量矢量场上表示论的自然推论。
$ISO(2)$ 小群的``平移"生成元在物理态上本征值为零,
但在 $A^\mu$ 上产生了 $\delta A_\mu = \partial_\mu\Omega$ 的规范变换。
}}
\end{center}

从表示论看, 取反对称导数实现了映射
\begin{equation}
    \underbrace{A^\mu}_{(\frac{1}{2},\frac{1}{2})}
    \xrightarrow{\;\partial_\mu A_\nu - \partial_\nu A_\mu\;}
    \underbrace{f_{\mu\nu}}_{(1,0)\oplus(0,1)}
\end{equation}
``投射掉"了 $j=0$ 分量, 只保留 $j=1$ 的两个横向自由度。

矢量势的场展开 (辐射规范 $A^0=0$, $\nabla\cdot\vec{A}=0$) 为
\begin{equation}
    A^\mu(x) = \frac{1}{(2\pi)^{3/2}} \sum_{\lambda=\pm 1}
    \int \frac{d^3k}{\sqrt{2k^0}}
    \left[\epsilon^\mu(\mathbf{k},\lambda)\,a(\mathbf{k},\lambda)\,e^{ik\cdot x}
    + \epsilon^{*\mu}(\mathbf{k},\lambda)\,a^\dagger(\mathbf{k},\lambda)\,e^{-ik\cdot x}\right]
\end{equation}
其中产生湮灭算符与 (\ref{eq:f-expansion}) 中的完全相同。
取反对称导数, 利用 $\partial_\mu e^{\pm ik\cdot x}=\pm ik_\mu e^{\pm ik\cdot x}$,
精确回到 (\ref{eq:f-expansion}): $\partial_\mu A_\nu - \partial_\nu A_\mu = f_{\mu\nu}$。


\topictitle{为什么需要矢量势}

如果 $f_{\mu\nu}$ 已经完整描述了自由光子, 为什么还需要 $A^\mu$?

\begin{enumerate}
    \item \textbf{局域相互作用}。由 2.5 节, 相互作用密度必须写成局域场在同一时空点的多项式。
    光子与带电粒子的最小耦合为 $\mathscr{H}_{\text{int}}=j^\mu A_\mu$ (维度 4, 可重整);
    若用 $f_{\mu\nu}$ 耦合则得到磁矩型 $j^{\mu\nu}f_{\mu\nu}$ (维度 5, 不可重整)。

    \item \textbf{规范不变性作为代价}。引入 $A^\mu$ 带来规范冗余, 需要规范固定和 Faddeev-Popov 程序处理,
    但换来的是可重整性和局域性。

    \item \textbf{Aharonov-Bohm 效应}。在量子理论中,
    $\exp(ie\oint A_\mu\,dx^\mu)$ 是可观测的——即使 $f_{\mu\nu}=0$ 的区域。
    这表明 $A_\mu$ 在量子理论中具有超越 $f_{\mu\nu}$ 的物理内容。
\end{enumerate}

总结: $(1,0)\oplus(0,1)$ 提供了自由光子的完整描述;
$(\frac{1}{2},\frac{1}{2})$ 是引入相互作用的必要工具。
```学上优雅, 但物理上不是必须的。
上面的推导只用到了三件事：电磁场的 boost 变换规则 (\ref{eq:boost-EB})、
复矢量的定义 $\vec{f}^{\,\pm}=\frac{1}{2}(\vec{E}\pm i\vec{B})$、
以及 Lorentz 代数 $\boldsymbol{A}=\frac{1}{2}(\boldsymbol{J}-i\boldsymbol{K})$。